\documentclass[10pt,a4paper]{amsart}
\usepackage{amsmath,amssymb,amsthm,graphicx}
\usepackage{verbatim}
\newenvironment{reference}{\comment}{\endcomment}
%\newenvironment{reference}{}{}
\newenvironment{slogan}{\comment}{\endcomment}
\newenvironment{history}{\comment}{\endcomment}

% For commutative diagrams we use Xy-pic
\usepackage[all]{xy}

% We use 2cell for 2-commutative diagrams.
\xyoption{2cell}
% Theorem environments.
%
\theoremstyle{plain}
\newtheorem{theorem}[subsection]{Theorem}
\newtheorem{proposition}[subsection]{Proposition}
\newtheorem{lemma}[subsection]{Lemma}

\theoremstyle{definition}
\newtheorem{definition}[subsection]{Definition}
\newtheorem{example}[subsection]{Example}
\newtheorem{exercise}[subsection]{Exercise}
\newtheorem{situation}[subsection]{Situation}

\theoremstyle{remark}
\newtheorem{remark}[subsection]{Remark}
\newtheorem{remarks}[subsection]{Remarks}

\numberwithin{equation}{subsection}

% Macros
%
\def\lim{\mathop{\mathrm{lim}}\nolimits}
\def\colim{\mathop{\mathrm{colim}}\nolimits}
\def\Spec{\mathop{\mathrm{Spec}}}
\def\Hom{\mathop{\mathrm{Hom}}\nolimits}
\def\Ext{\mathop{\mathrm{Ext}}\nolimits}
\def\SheafHom{\mathop{\mathcal{H}\!\mathit{om}}\nolimits}
\def\SheafExt{\mathop{\mathcal{E}\!\mathit{xt}}\nolimits}
\def\Sch{\mathit{Sch}}
\def\Mor{\mathop{\mathrm{Mor}}\nolimits}
\def\Ob{\mathop{\mathrm{Ob}}\nolimits}
\def\Sh{\mathop{\mathit{Sh}}\nolimits}
\def\NL{\mathop{N\!L}\nolimits}
\def\CH{\mathop{\mathrm{CH}}\nolimits}
\def\proetale{{pro\text{-}\acute{e}tale}}
\def\etale{{\acute{e}tale}}
\def\QCoh{\mathit{QCoh}}
\def\Ker{\mathop{\mathrm{Ker}}}
\def\Im{\mathop{\mathrm{Im}}}
\def\Coker{\mathop{\mathrm{Coker}}}
\def\Coim{\mathop{\mathrm{Coim}}}

% Boxtimes
%
\DeclareMathSymbol{\boxtimes}{\mathbin}{AMSa}{"02}

%
% Macros for moduli stacks/spaces
%
\def\QCohstack{\mathcal{QC}\!\mathit{oh}}
\def\Cohstack{\mathcal{C}\!\mathit{oh}}
\def\Spacesstack{\mathcal{S}\!\mathit{paces}}
\def\Quotfunctor{\mathrm{Quot}}
\def\Hilbfunctor{\mathrm{Hilb}}
\def\Curvesstack{\mathcal{C}\!\mathit{urves}}
\def\Polarizedstack{\mathcal{P}\!\mathit{olarized}}
\def\Complexesstack{\mathcal{C}\!\mathit{omplexes}}
% \Pic is the operator that assigns to X its picard group, usage \Pic(X)
% \Picardstack_{X/B} denotes the Picard stack of X over B
% \Picardfunctor_{X/B} denotes the Picard functor of X over B
\def\Pic{\mathop{\mathrm{Pic}}\nolimits}
\def\Picardstack{\mathcal{P}\!\mathit{ic}}
\def\Picardfunctor{\mathrm{Pic}}
\def\Deformationcategory{\mathcal{D}\!\mathit{ef}}
\usepackage{cleveref}
\crefname{originalequation}{original equation}{original equations}
\Crefname{originalequation}{Original equation}{Original equations}
\makeatletter
\newlabel{equation-multiplicity}{{0BI6}{source}{}{sourcepointer-1}{}}
\newlabel{equation-multiplicity@cref}{{[originalequation][0][]0BI6}{[1][0][]source}{}{}{}}
\makeatother
\crefname{originallemma}{original lemma}{original lemmas}
\Crefname{originallemma}{Original lemma}{Original lemmas}
\makeatletter
\newlabel{etale-lemma-strict-normal-crossings}{{0BIA}{source}{}{sourcepointer-2}{}}
\newlabel{etale-lemma-strict-normal-crossings@cref}{{[originallemma][0][]0BIA}{[1][0][]source}{}{}{}}
\makeatother
\crefname{originalreference}{original reference}{original references}
\Crefname{originalreference}{Original reference}{Original references}
\makeatletter
\newlabel{lemma-blowup}{{0AGQ}{source}{}{sourcepointer-3}{}}
\newlabel{lemma-blowup@cref}{{[originalreference][0][]0AGQ}{[1][0][]source}{}{}{}}
\makeatother
\crefname{originalreference}{original reference}{original references}
\Crefname{originalreference}{Original reference}{Original references}
\makeatletter
\newlabel{lemma-blowup-curves}{{0BI8}{source}{}{sourcepointer-4}{}}
\newlabel{lemma-blowup-curves@cref}{{[originalreference][0][]0BI8}{[1][0][]source}{}{}{}}
\makeatother
\crefname{originalreference}{original reference}{original references}
\Crefname{originalreference}{Original reference}{Original references}
\makeatletter
\newlabel{lemma-blowup-nonsingular-curves-meeting-at-point}{{0BI7}{source}{}{sourcepointer-5}{}}
\newlabel{lemma-blowup-nonsingular-curves-meeting-at-point@cref}{{[originalreference][0][]0BI7}{[1][0][]source}{}{}{}}
\makeatother
\crefname{originalreference}{original reference}{original references}
\Crefname{originalreference}{Original reference}{Original references}
\makeatletter
\newlabel{lemma-resolve-curve}{{0BI4}{source}{}{sourcepointer-6}{}}
\newlabel{lemma-resolve-curve@cref}{{[originalreference][0][]0BI4}{[1][0][]source}{}{}{}}
\makeatother
\crefname{originalreference}{original reference}{original references}
\Crefname{originalreference}{Original reference}{Original references}
\makeatletter
\newlabel{lemma-turn-into-effective-Cartier}{{0BIB}{source}{}{sourcepointer-7}{}}
\newlabel{lemma-turn-into-effective-Cartier@cref}{{[originalreference][0][]0BIB}{[1][0][]source}{}{}{}}
\makeatother
\title{Embedded resolution of a proper closed subscheme of a regular surface by point blowups}
\author{The Stacks Project Authors}
\date{stacks-project:resolve.tex:local-snapshot-20261012; GNU FDL 1.2 or later}
\begin{document}
\maketitle
\noindent\textit{Editorial note.} All delivered mathematical text is the author's own, copied from the pinned Stacks Project chapter file resolve.tex; no statement or proof step was rewritten, reordered or removed.\par\medskip
\noindent\textit{Editorial note.} The counted proof is the author's complete argument: the reduction to an effective Cartier divisor with regular components using the description of strict transforms and the exceptional divisor, the descending induction on the maximal intersection multiplicity of two components, and the separation of the points where more than two components meet leading to the strict normal crossings condition.\par\medskip
\noindent\textit{Editorial note.} The label lemma-embedded-resolution names the counted statement, and the proof environment immediately following it in the source is delivered verbatim as its proof.\par\medskip
\noindent\textit{Editorial note.} The counted statement is carried by the author's own lemma environment in the chapter source.\par\medskip
\noindent\textit{Editorial note.} Pre-open gate: python operations/agent-authoring/tools/precard\_gate.py --source resolve.tex --label lemma-embedded-resolution --start 4447 --end 4501 --local-id 28080 returned RESULT=OK (L28080 free, occupancy FREE, public interval clear). Pre-assembly tag self-check: canonical tag resolved and not present in reference/public-items.jsonl.\par\medskip
\noindent\textit{Editorial note.} The statement and proof are copied from the pinned chapter file; the cited inputs are kept as external pointers with their canonical Stacks tags.\par\medskip
\noindent\textit{Allowed standard prerequisites (editorial).}\begin{itemize}
\item Behaviour of the inverse image of the subscheme under a point blowup. Same chapter, the lemma describing the blowup of a regular surface at a closed point and the description of the strict transform of a curve under such a blowup. Supplies the facts that the exceptional divisor is a regular curve and that the strict transform of a component is either the component itself or its blowup at the point, so no new singular one dimensional components appear.
\item Making the subscheme an effective Cartier divisor with regular components. Same chapter, the lemnas turning the subscheme of the surface into an effective Cartier divisor and making its components regular by blowups in closed points. Supplies the reduction to the case where the subscheme is an effective Cartier divisor with regular irreducible components.
\item Multiplicity invariant of the intersection of two regular curves and its behaviour under blowup. Same chapter, the definition of the multiplicity of the intersection of two curves and the lemma computing its behaviour under a blowup of the intersection points. Supplies the invariant used for the descending induction, together with the fact that the new invariants against the exceptional divisor are one.
\item Criterion for a divisor with regular components to be strict normal crossings. Etale Morphisms, the lemma characterizing a strict normal crossings divisor through the local shape of the components at every point. Supplies the final recognition of the divisor as a strict normal crossings divisor.
\end{itemize}\medskip
\ifdefined\hypertarget\hypertarget{sourcepointer-1}{}\fi
\noindent\textit{Source reference.} Original equation 0BI6: sources/stacks-resolve/pinned-original-source.tex, source lines 4333--4333. Prerequisite pointer invoked at proof line(s) 4480; it names a standard fact of the same chapter that the counted argument uses, and every step of the counted proof built on it is retained verbatim in the delivered body.\par\medskip
\ifdefined\hypertarget\hypertarget{sourcepointer-2}{}\fi
\noindent\textit{Source reference.} Original lemma 0BIA: sources/stacks-chow/tags.txt, source lines 14483--14483. Prerequisite pointer invoked at proof line(s) 4500; it names a standard fact of a sibling chapter that the counted argument uses, and every step of the counted proof built on it is retained verbatim in the delivered body.\par\medskip
\ifdefined\hypertarget\hypertarget{sourcepointer-3}{}\fi
\noindent\textit{Source reference.} Original reference 0AGQ: sources/stacks-resolve/pinned-original-source.tex, source lines 666--666. Prerequisite pointer invoked at proof line(s) 4465; it names a standard fact of the same chapter that the counted argument uses, and every step of the counted proof built on it is retained verbatim in the delivered body.\par\medskip
\ifdefined\hypertarget\hypertarget{sourcepointer-4}{}\fi
\noindent\textit{Source reference.} Original reference 0BI8: sources/stacks-resolve/pinned-original-source.tex, source lines 4395--4395. Prerequisite pointer invoked at proof line(s) 4469; it names a standard fact of the same chapter that the counted argument uses, and every step of the counted proof built on it is retained verbatim in the delivered body.\par\medskip
\ifdefined\hypertarget\hypertarget{sourcepointer-5}{}\fi
\noindent\textit{Source reference.} Original reference 0BI7: sources/stacks-resolve/pinned-original-source.tex, source lines 4346--4346. Prerequisite pointer invoked at proof line(s) 4482; it names a standard fact of the same chapter that the counted argument uses, and every step of the counted proof built on it is retained verbatim in the delivered body.\par\medskip
\ifdefined\hypertarget\hypertarget{sourcepointer-6}{}\fi
\noindent\textit{Source reference.} Original reference 0BI4: sources/stacks-resolve/pinned-original-source.tex, source lines 4248--4248. Prerequisite pointer invoked at proof line(s) 4452; it names a standard fact of the same chapter that the counted argument uses, and every step of the counted proof built on it is retained verbatim in the delivered body.\par\medskip
\ifdefined\hypertarget\hypertarget{sourcepointer-7}{}\fi
\noindent\textit{Source reference.} Original reference 0BIB: sources/stacks-resolve/pinned-original-source.tex, source lines 4424--4424. Prerequisite pointer invoked at proof line(s) 4469; it names a standard fact of the same chapter that the counted argument uses, and every step of the counted proof built on it is retained verbatim in the delivered body.\par\medskip
\begin{lemma}
\label{lemma-embedded-resolution}
Let $X$ be a regular scheme of dimension $2$. Let $Z \subset X$
be a proper closed subscheme such that every irreducible component
$Y \subset Z$ of dimension $1$ satisfies the equivalent conditions of
Lemma \ref{lemma-resolve-curve}. Then there exists a sequence
$$
X_n \to \ldots \to X_1 \to X
$$
of blowups in closed points such that the inverse image $Z_n$ of $Z$
in $X_n$ is an effective Cartier divisor supported on a strict normal crossings
divisor.
\end{lemma}

\begin{proof}
Let $X' \to X$ be a blowup in a closed point $p$. Then the inverse image
$Z' \subset X'$ of $Z$ is supported on the strict transform of $Z$ and
the exceptional divisor. The exceptional divisor is a regular curve
(Lemma \ref{lemma-blowup}) and the strict transform $Y'$ of each irreducible
component $Y$ is either equal to $Y$ or the blowup of $Y$ at $p$.
Thus in this process we do not produce additional singular components
of dimension $1$. Thus it follows from
Lemmas \ref{lemma-turn-into-effective-Cartier} and \ref{lemma-blowup-curves}
that we may assume $Z$ is an effective Cartier divisor and
that all irreducible components $Y$ of $Z$ are regular.
(Of course we cannot assume the irreducible components are
pairwise disjoint because in each blowup of a point of $Z$
we add a new irreducible component to $Z$, namely the exceptional divisor.)

\medskip\noindent
Assume $Z$ is an effective Cartier divisor whose irreducible components
$Y_i$ are regular. For every $i \not = j$
and $p \in Y_i \cap Y_j$ we have the invariant
$m_p(Y_i \cap Y_j)$ (\ref{equation-multiplicity}).
If the maximum of these numbers is $> 1$, then we can decrease
it (Lemma \ref{lemma-blowup-nonsingular-curves-meeting-at-point})
by blowing up in all the points $p$ where the maximum is attained
(note that the ``new'' invariants $m_{q_i}(Y'_i \cap E)$ are always $1$).
If the maximum is $1$ then, if $p \in Y_1 \cap \ldots \cap Y_r$
for some $r > 2$ and not any of the others (for example), then after
blowing up $p$ we see that $Y'_1, \ldots, Y'_r$ do not meet in points
above $p$ and $m_{q_i}(Y'_i, E) = 1$ where $Y'_i \cap E = \{q_i\}$.
Thus continuing to blowup points where more than $3$
of the components of $Z$ meet, we reach the situation where
for every closed point $p \in X$ there is either
(a) no curves $Y_i$ passing through $p$,
(b) exactly one curve $Y_i$ passing through $p$ and $\mathcal{O}_{Y_i, p}$
is regular, or (c) exactly two curves $Y_i$, $Y_j$ passing through
$p$, the local rings $\mathcal{O}_{Y_i, p}$, $\mathcal{O}_{Y_j, p}$
are regular and $m_p(Y_i \cap Y_j) = 1$.
This means that $\sum Y_i$ is a strict normal crossings
divisor on the regular surface $X$, see
\'Etale Morphisms, Lemma
\ref{etale-lemma-strict-normal-crossings}.
\end{proof}

\end{document}
